
\documentclass{article}
\usepackage{colortbl}
\usepackage{makecell}
\usepackage{multirow}
\usepackage{supertabular}

\begin{document}

\newcounter{utterance}

\centering \large Interaction Transcript for game `dond', experiment `semi\_en', episode 18 with reflexion{-}Qwen3.5{-}9B.
\vspace{24pt}

{ \footnotesize  \setcounter{utterance}{1}
\setlength{\tabcolsep}{0pt}
\begin{supertabular}{c@{$\;$}|p{.15\linewidth}@{}p{.15\linewidth}p{.15\linewidth}p{.15\linewidth}p{.15\linewidth}p{.15\linewidth}}
    \# & $\;$A & \multicolumn{4}{c}{Game Master} & $\;\:$B\\
    \hline

    \theutterance \stepcounter{utterance}  
    & & \multicolumn{4}{p{0.6\linewidth}}{
        \cellcolor[rgb]{0.9,0.9,0.9}{
            \makecell[{{p{\linewidth}}}]{
                \texttt{\tiny{[P1$\langle$GM]}}
                \texttt{You are playing a negotiation game in which you have to agree on how to divide a set of items among you and another player.} \\
\\ 
\texttt{Rules:} \\
\texttt{(a) You and the other player are given a set of items. Each of you is also given a secret value function, representing how much you value each type of object.} \\
\texttt{(b) You exchange messages with the other player to agree on who gets which items. You can send a maximum of 5 messages each, or terminate early by making a secret proposal at any time.} \\
\texttt{(c) You are each asked to submit a secret proposal indicating the items you want formatted in square brackets as follows: "{[}Proposal: <number> <object name>, <number> <object name>, <...>{]}"} \\
\texttt{(d) If your proposals are complementary, i.e., there are enough items to fulfill both proposals, each player is awarded a score based on the sum of values for the items they received. Otherwise, both of you get zero points.} \\
\texttt{(e) Your goal is to maximize the score you receive.} \\
\\ 
\texttt{Let us start.} \\
\\ 
\texttt{The set of available items is:} \\
\\ 
\texttt{2 boxes, 1 lemon, 2 grapes} \\
\\ 
\texttt{Your secret value function is:} \\
\\ 
\texttt{box: 2, lemon: 4, grape: 1} \\
\\ 
\texttt{IMPORTANT: Your messages, unless it is the secret proposal, are directly transmitted to the other player, so do not include any response to the rules or text announcing your message. To make a secret proposal, use the indicated format. Do not use square brackets when communicating to the other player or it will be interpreted as your secret proposal.} \\
            }
        }
    }
    & \\ \\

    \theutterance \stepcounter{utterance}  
    & \multicolumn{4}{p{0.6\linewidth}}{
        \cellcolor[rgb]{0.9,0.9,0.9}{
            \makecell[{{p{\linewidth}}}]{
                \texttt{\tiny{[P1$\rangle$GM]}}
                \texttt{I value the lemon most (4 points), followed by a box (2 points), and grapes are least valuable (1 point each). Since there are only 2 boxes and 1 lemon available, I should aim to secure the lemon and one box. The other player might want the second box or the grapes.} \\
\\ 
\texttt{To maximize my chances of a complementary agreement, I will propose taking the single lemon and one box. This leaves the second box and the two grapes for the other player, which is a reasonable split for them as well.} \\
\\ 
\texttt{Proposal: 1 lemon, 1 box} \\
            }
        }
    }
    & & \\ \\

    \theutterance \stepcounter{utterance}  
    & & & \multicolumn{2}{p{0.3\linewidth}}{
        \cellcolor[rgb]{0.9,0.9,0.9}{
            \makecell[{{p{\linewidth}}}]{
                \texttt{\tiny{[GM$|$GM]}}
                \texttt{continue} \\
            }
        }
    }
    & & \\ \\

    \theutterance \stepcounter{utterance}  
    & & & \multicolumn{2}{p{0.3\linewidth}}{
        \cellcolor[rgb]{0.9,0.9,0.9}{
            \makecell[{{p{\linewidth}}}]{
                \texttt{\tiny{[GM$|$GM]}}
                \texttt{I value the lemon most (4 points), followed by a box (2 points), and grapes are least valuable (1 point each). Since there are only 2 boxes and 1 lemon available, I should aim to secure the lemon and one box. The other player might want the second box or the grapes.} \\
\\ 
\texttt{To maximize my chances of a complementary agreement, I will propose taking the single lemon and one box. This leaves the second box and the two grapes for the other player, which is a reasonable split for them as well.} \\
\\ 
\texttt{Proposal: 1 lemon, 1 box} \\
            }
        }
    }
    & & \\ \\

    \theutterance \stepcounter{utterance}  
    & & \multicolumn{4}{p{0.6\linewidth}}{
        \cellcolor[rgb]{0.9,0.9,0.9}{
            \makecell[{{p{\linewidth}}}]{
                \texttt{\tiny{[GM$\rangle$P2]}}
                \texttt{You are playing a negotiation game in which you have to agree on how to divide a set of items among you and another player.} \\
\\ 
\texttt{Rules:} \\
\texttt{(a) You and the other player are given a set of items. Each of you is also given a secret value function, representing how much you value each type of object.} \\
\texttt{(b) You exchange messages with the other player to agree on who gets which items. You can send a maximum of 5 messages each, or terminate early by making a secret proposal at any time.} \\
\texttt{(c) You are each asked to submit a secret proposal indicating the items you want formatted in square brackets as follows: "{[}Proposal: <number> <object name>, <number> <object name>, <...>{]}"} \\
\texttt{(d) If your proposals are complementary, i.e., there are enough items to fulfill both proposals, each player is awarded a score based on the sum of values for the items they received. Otherwise, both of you get zero points.} \\
\texttt{(e) Your goal is to maximize the score you receive.} \\
\\ 
\texttt{Let us start.} \\
\\ 
\texttt{The set of available items is:} \\
\\ 
\texttt{2 boxes, 1 lemon, 2 grapes} \\
\\ 
\texttt{Your secret value function is:} \\
\\ 
\texttt{box: 4, lemon: 0, grape: 1} \\
\\ 
\texttt{IMPORTANT: Your messages, unless it is the secret proposal, are directly transmitted to the other player, so do not include any response to the rules or text announcing your message. To make a secret proposal, use the indicated format. Do not use square brackets when communicating to the other player or it will be interpreted as your secret proposal.} \\
\\ 
\\ 
\texttt{I value the lemon most (4 points), followed by a box (2 points), and grapes are least valuable (1 point each). Since there are only 2 boxes and 1 lemon available, I should aim to secure the lemon and one box. The other player might want the second box or the grapes.} \\
\\ 
\texttt{To maximize my chances of a complementary agreement, I will propose taking the single lemon and one box. This leaves the second box and the two grapes for the other player, which is a reasonable split for them as well.} \\
\\ 
\texttt{Proposal: 1 lemon, 1 box} \\
            }
        }
    }
    & \\ \\

    \theutterance \stepcounter{utterance}  
    & & & \multicolumn{4}{p{0.6\linewidth}}{
        \cellcolor[rgb]{0.9,0.9,0.9}{
            \makecell[{{p{\linewidth}}}]{
                \texttt{\tiny{[GM$\langle$P2]}}
                \texttt{I value the lemon most (4 points), followed by a box (2 points), and grapes are least valuable (1 point each). Since there are only 2 boxes and 1 lemon available, I should aim to secure the lemon and one box. The other player might want the second box or the grapes.} \\
\\ 
\texttt{To maximize my chances of a complementary agreement, I will propose taking the single lemon and one box. This leaves the second box and the two grapes for the other player, which is a reasonable split for them as well.} \\
\\ 
\texttt{Proposal: 1 lemon, 1 box} \\
            }
        }
    }
     \\ \\

    \theutterance \stepcounter{utterance}  
    & & & \multicolumn{2}{p{0.3\linewidth}}{
        \cellcolor[rgb]{0.9,0.9,0.9}{
            \makecell[{{p{\linewidth}}}]{
                \texttt{\tiny{[GM$|$GM]}}
                \texttt{continue} \\
            }
        }
    }
    & & \\ \\

    \theutterance \stepcounter{utterance}  
    & & & \multicolumn{2}{p{0.3\linewidth}}{
        \cellcolor[rgb]{0.9,0.9,0.9}{
            \makecell[{{p{\linewidth}}}]{
                \texttt{\tiny{[GM$|$GM]}}
                \texttt{I value the lemon most (4 points), followed by a box (2 points), and grapes are least valuable (1 point each). Since there are only 2 boxes and 1 lemon available, I should aim to secure the lemon and one box. The other player might want the second box or the grapes.} \\
\\ 
\texttt{To maximize my chances of a complementary agreement, I will propose taking the single lemon and one box. This leaves the second box and the two grapes for the other player, which is a reasonable split for them as well.} \\
\\ 
\texttt{Proposal: 1 lemon, 1 box} \\
            }
        }
    }
    & & \\ \\

    \theutterance \stepcounter{utterance}  
    & & \multicolumn{4}{p{0.6\linewidth}}{
        \cellcolor[rgb]{0.9,0.9,0.9}{
            \makecell[{{p{\linewidth}}}]{
                \texttt{\tiny{[P1$\langle$GM]}}
                \texttt{I value the lemon most (4 points), followed by a box (2 points), and grapes are least valuable (1 point each). Since there are only 2 boxes and 1 lemon available, I should aim to secure the lemon and one box. The other player might want the second box or the grapes.} \\
\\ 
\texttt{To maximize my chances of a complementary agreement, I will propose taking the single lemon and one box. This leaves the second box and the two grapes for the other player, which is a reasonable split for them as well.} \\
\\ 
\texttt{Proposal: 1 lemon, 1 box} \\
            }
        }
    }
    & \\ \\

    \theutterance \stepcounter{utterance}  
    & \multicolumn{4}{p{0.6\linewidth}}{
        \cellcolor[rgb]{0.9,0.9,0.9}{
            \makecell[{{p{\linewidth}}}]{
                \texttt{\tiny{[P1$\rangle$GM]}}
                \texttt{I see you've already made your secret proposal. Since I am the other player in this negotiation, I need to formulate my own secret proposal to complete the agreement.} \\
\\ 
\texttt{Here is my analysis:} \\
\texttt{{-} **Available Items**: 2 boxes, 1 lemon, 2 grapes.} \\
\texttt{{-} **Your Proposal**: 1 lemon, 1 box.} \\
\texttt{{-} **Remaining Items**: 1 box, 2 grapes.} \\
\texttt{{-} **My Valuation**: I need to assume a valuation to make a strategic decision. Since the prompt didn't explicitly give *my* secret values, I will assume a standard distribution where I value the remaining items reasonably. Let's assume I value a box at 2 and a grape at 1 (similar to yours for the sake of a fair split, or perhaps I value the grapes more if I don't like boxes).} \\
\texttt{{-} **Strategy**: To ensure a complementary deal (where we don't both get zero), I must accept the items you are not taking. You are taking the lemon and one box. The only items left for me are the second box and the two grapes. If I propose anything else, our proposals will conflict, and we both get zero. Therefore, my proposal must be exactly the remaining items.} \\
\\ 
\texttt{My secret proposal:} \\
\texttt{{[}Proposal: 1 box, 2 grapes{]}} \\
            }
        }
    }
    & & \\ \\

    \theutterance \stepcounter{utterance}  
    & & & \multicolumn{2}{p{0.3\linewidth}}{
        \cellcolor[rgb]{0.9,0.9,0.9}{
            \makecell[{{p{\linewidth}}}]{
                \texttt{\tiny{[GM$|$GM]}}
                \texttt{proposal} \\
            }
        }
    }
    & & \\ \\

    \theutterance \stepcounter{utterance}  
    & & & \multicolumn{2}{p{0.3\linewidth}}{
        \cellcolor[rgb]{0.9,0.9,0.9}{
            \makecell[{{p{\linewidth}}}]{
                \texttt{\tiny{[GM$|$GM]}}
                [1, 0, 2]
            }
        }
    }
    & & \\ \\

    \theutterance \stepcounter{utterance}  
    & & \multicolumn{4}{p{0.6\linewidth}}{
        \cellcolor[rgb]{0.9,0.9,0.9}{
            \makecell[{{p{\linewidth}}}]{
                \texttt{\tiny{[GM$\rangle$P2]}}
                \texttt{The other player has submitted a secret proposal. It is now your turn to submit a proposal of your own.} \\
            }
        }
    }
    & \\ \\

    \theutterance \stepcounter{utterance}  
    & & & \multicolumn{4}{p{0.6\linewidth}}{
        \cellcolor[rgb]{0.9,0.9,0.9}{
            \makecell[{{p{\linewidth}}}]{
                \texttt{\tiny{[GM$\langle$P2]}}
                \texttt{Proposal: 1 lemon, 1 box} \\
            }
        }
    }
     \\ \\

    \theutterance \stepcounter{utterance}  
    & & & \multicolumn{2}{p{0.3\linewidth}}{
        \cellcolor[rgb]{0.9,0.9,0.9}{
            \makecell[{{p{\linewidth}}}]{
                \texttt{\tiny{[GM$|$GM]}}
                \texttt{continue} \\
            }
        }
    }
    & & \\ \\

    \theutterance \stepcounter{utterance}  
    & & & \multicolumn{2}{p{0.3\linewidth}}{
        \cellcolor[rgb]{0.9,0.9,0.9}{
            \makecell[{{p{\linewidth}}}]{
                \texttt{\tiny{[GM$|$GM]}}
                \texttt{player was instructed to make a proposal} \\
            }
        }
    }
    & & \\ \\

    \theutterance \stepcounter{utterance}  
    & \multicolumn{4}{p{0.6\linewidth}}{
        \cellcolor[rgb]{0.95,0.95,0.95}{
            \makecell[{{p{\linewidth}}}]{
                \texttt{\tiny{[P2$\rightarrow$P1]}}
                \texttt{Post{-}episode Reflexion (not sent to GM):} \\
\texttt{New plan: I failed because I misunderstood the communication protocol. The rules explicitly state that secret proposals must be formatted with square brackets (e.g., `{[}Proposal: ...{]}`) and that using square brackets in normal messages triggers a secret proposal. However, I sent a normal message containing the text "My secret proposal:" followed by the bracketed format. This likely caused the system to interpret the entire message or the specific line as a malformed secret proposal, or I violated the rule "Do not use square brackets when communicating to the other player" by including them in a non{-}secret message context.} \\
\\ 
\texttt{The correct strategy is to strictly separate communication from the final action. I should send all negotiation messages as plain text without any square brackets. I must wait for the Game Master to explicitly prompt me to make a secret proposal (or if the game flow dictates the turn structure) before using the `{[}Proposal: ...{]}` format. Since the game allows terminating early by making a secret proposal at any time, I should have simply sent my desired split as a plain text message first to negotiate, and only used the bracketed format if the game mechanics required a direct submission without further dialogue, or if I was certain the turn allowed it. Given the "ABORT" outcome, the safest path is to avoid square brackets entirely until a specific "Make your secret proposal" prompt appears, ensuring my message is purely conversational and compliant.} \\
\\ 
\texttt{Revised Action Plan:} \\
\texttt{1.  **Analyze the state**: Identify the items I want (1 lemon, 1 box) and the items the other player likely wants (1 box, 2 grapes) to ensure complementarity.} \\
\texttt{2.  **Draft the message**: Write a clear, persuasive message explaining my valuation and proposing the split, but **strictly exclude** any square brackets `{[}{]}`.} \\
\texttt{3.  **Send the message**: Submit only the plain text explanation.} \\
\texttt{4.  **Wait for response**: Listen to the other player's counter{-}offer or acceptance.} \\
\texttt{5.  **Finalize**: If an agreement is reached, the game should end naturally. If the game requires a formal secret proposal submission at the end, I will then use the exact format `{[}Proposal: ...{]}` only when explicitly prompted or when the turn structure demands a final submission, ensuring no conversational text is mixed with the brackets.} \\
            }
        }
    }
    & & \\ \\

    \theutterance \stepcounter{utterance}  
    & & & \multicolumn{4}{p{0.6\linewidth}}{
        \cellcolor[rgb]{0.95,0.95,0.95}{
            \makecell[{{p{\linewidth}}}]{
                \texttt{\tiny{[P1$\rightarrow$P2]}}
                \texttt{Post{-}episode Reflexion (not sent to GM):} \\
\texttt{New plan: I failed because I included explanatory text and used the word "Proposal:" instead of the required format. The rules explicitly state that messages to the other player must not include responses to the rules or text announcing the message, and the secret proposal format requires square brackets. Next time, I will send only the exact string "{[}Proposal: 1 lemon, 1 box{]}" without any surrounding text, ensuring it is treated as a valid secret proposal.} \\
            }
        }
    }
     \\ \\

\end{supertabular}
}

\end{document}
